%% Nature Computational Science submission
%% Draft — [DATE]
\documentclass[pdflatex,sn-nature]{sn-jnl}

\usepackage{graphicx}
\usepackage{multirow}
\usepackage{amsmath,amssymb,amsfonts}
\usepackage{amsthm}
\usepackage[title]{appendix}
\usepackage{xcolor}
\usepackage{booktabs}
\usepackage{array}
\usepackage{tabularx}

\raggedbottom

\begin{document}

%%==============================================================%%
%% Title
%%==============================================================%%
\title[Temporal multimodal ALL prediction]{%
  Temporally-structured multimodal prediction across six clinical
  targets in paediatric acute lymphoblastic leukaemia}

%%==============================================================%%
%% Authors
%%==============================================================%%
\author*[1,2]{\fnm{xxx} \sur{xx}}\email{xxx}

\author[1]{\fnm{[Author 2 Given]} \sur{[Author 2 Family]}}

\author[1]{\fnm{[Author 3 Given]} \sur{[Author 3 Family]}}

\affil*[1]{\orgdiv{[Department]},
           \orgname{[HOSPITAL NAME / University]},
           \orgaddress{\city{[City]}, \country{China}}}

\affil[2]{\orgdiv{School of Computing and Information Systems},
           \orgname{The University of Melbourne},
           \orgaddress{\city{Melbourne}, \country{Australia}}}

%%==============================================================%%
%% Abstract
%%==============================================================%%
\abstract{%
Paediatric acute lymphoblastic leukaemia (ALL) is the most common
childhood malignancy, yet 15--20\% of patients relapse after
apparently successful treatment.
Existing risk-stratification approaches operate on single-timepoint
snapshots and have not been systematically evaluated across the full
clinical decision timeline nor validated across multiple independent
centres.
Here we present a temporally-structured multimodal prediction framework
evaluated on six clinical targets simultaneously: four classification
tasks (Final Risk, D19MRD, D46MRD, Relapse) and two survival endpoints
(event-free survival, EFS; overall survival, OS), whose feature sets
are strictly aligned with the clinical data-acquisition timeline from
hospital admission to Day~46.
Phase~1 evaluates four large language models (LLMs) in a zero-shot
setting; Phase~2 trains decision trees, random forests, and random
survival forests on a temporally partitioned case database.
In the internal cohort, Phase~2 achieves AUROC~0.949 for Final Risk
and concordance index~0.746 for EFS.
A cosine similarity retrieval-augmented generation (RAG) belief
revision improves D46MRD AUROC from 0.655 to 0.881 ($+0.226$) in
the largest training regime.
External validation at an independent site confirms robust Final Risk
performance (RF AUROC~0.991) and LLM OS concordance indices
of 0.806--0.859 without retraining.
This framework demonstrates that temporally stratified multimodal
learning, aligned with clinical workflows, enables consistent
multi-target prediction that generalises to independent cohorts.
}

\keywords{paediatric acute lymphoblastic leukaemia, temporal feature
design, large language models, risk stratification, survival analysis,
retrieval-augmented generation, external validation}

\maketitle

%%==============================================================%%
\section{Introduction}\label{sec:intro}
%%==============================================================%%

Acute lymphoblastic leukaemia (ALL) is the most common childhood
malignancy, accounting for approximately 25\% of all paediatric
cancers~\cite{[CITE:Hunger2015]}.
Risk-adapted chemotherapy has elevated five-year event-free survival
(EFS) rates above 90\% in high-income settings, yet the 15--20\% of
patients who relapse face substantially worse outcomes, with long-term
survival dropping to 30--50\% upon re-induction~\cite{[CITE:Bhojwani2013],[CITE:Pui2008]}.
The majority of relapses occur after apparent clinical remission,
making prospective identification of high-risk patients a central and
unresolved challenge.

Current standard-of-care risk stratification—Initial Risk at
diagnosis; Final Risk after Day~46 induction assessment—relies on
pre-specified clinical cut-offs for age, white blood cell count,
cytogenetics, immunophenotype, and minimal residual disease (MRD).
These categorical rules do not fully exploit the rich multimodal
information accumulated progressively during treatment, and do not
yield personalised probability estimates for future relapse or
survival.
Moreover, most published machine learning models operate on a fixed,
single-timepoint feature snapshot at diagnosis, ignoring the structured
temporal ordering in which data actually become clinically
available~\cite{[CITE:Steyerberg2019]}.
A model trained on pooled features from all timepoints may implicitly
rely on information unavailable at the point of decision-making,
producing optimistic but non-deployable estimates.

A second limitation is the absence of unified six-target evaluation:
classification tasks (Final Risk, D19MRD, D46MRD, Relapse) and survival
endpoints (EFS, OS) are typically studied in isolation, obscuring
how a single feature-acquisition workflow supports sequential
predictions across the entire induction period.
Third, the emergence of large language models (LLMs) with strong
zero-shot reasoning opens a complementary, training-free avenue for
clinical prediction~\cite{[CITE:Phaterpekar2026]}, yet systematic
evaluation on paediatric oncology outcome prediction remains absent.
Fourth, most single-centre machine learning studies do not include
external validation~\cite{[CITE:Martinez2025]}, limiting confidence
in cross-institution generalisability.

Here we address these gaps with a \emph{temporally-structured}
prediction framework applied to a single-centre training cohort and
validated at independent external sites.
Our contributions are: (\textit{i})~a twelve-level feature hierarchy
$\mathcal{F}_1 \subseteq \cdots \subseteq \mathcal{F}_{12}$ whose
boundaries align with the clinical data-acquisition timeline (Day~0
through Day~46+), applied jointly to six prediction targets;
(\textit{ii})~Phase~1, a systematic zero-shot evaluation of four LLMs
across all feature levels and all six targets, establishing
training-free baselines; (\textit{iii})~Phase~2, a decision tree,
random forest, and random survival forest system trained on a dynamic
case database partitioned by enrolment year, validating performance
under three temporal splits; (\textit{iv})~a bone marrow morphology
pipeline converting microscopy images to structured text via
SAM-based segmentation and BiomedCLIP clustering; (\textit{v})~a
cosine RAG belief revision that applies similarity-weighted updates to
RF predicted probabilities and generates interpretable LLM clinical
narratives; and (\textit{vi})~multi-site external validation at an
independent centre that received no training data.

%%==============================================================%%
\section{Results}\label{sec:results}
%%==============================================================%%

\subsection{Study cohorts and temporally-structured feature design}
\label{sec:cohort}

The internal training cohort comprised paediatric ALL patients treated
at [HOSPITAL NAME] between 2015 and 2019 (Table~\ref{tab:cohort};
$n=301$).
Prediction targets span the treatment timeline: Day~19 MRD positivity
(D19MRD $\geq 1\%$), Day~46 MRD positivity (D46MRD $\geq 0.01\%$),
final risk classification (Final Risk: LR vs IR/HR; 159 LR, 142 IR/HR),
and relapse occurrence (44 of 301, 14.6\%).
EFS events (relapse, induction failure, or death; 75 of 301, 24.9\%)
and OS events (death; 24 of 301, 8.0\%) together extend the prediction
targets to six.

Two external validation cohorts were recruited from independent
institutions.
External Site~1 (Lianyungang First People's Hospital) contributed
41 patients diagnosed after 2019, of whom 29 had complete labels for
the four classification tasks (Table~\ref{tab:cohort}).
External Site~2 ([EXTERNAL HOSPITAL 2 NAME]) data collection is
in progress; a placeholder row is included in Table~\ref{tab:cohort}
and the corresponding result cells are left empty.
For external validation, all 301 internal patients were used as the
training set and no external data were used in training.

Feature sets $\mathcal{F}_1 \subseteq \cdots \subseteq \mathcal{F}_{12}$
(Table~\ref{tab:expdesign}) are strictly nested: each level adds only
features available at the corresponding clinical timepoint.
$\mathcal{F}_1$ contains WBC count and age (Day~0);
$\mathcal{F}_8$ additionally includes D19MRD and D46MRD values (Day~46).
The feature hierarchy applies identically to all six prediction
targets, with the distinction that classification tasks (targets 1--4)
use AUROC as the primary metric, while survival targets (targets 5--6:
EFS, OS) use Harrell's C-index and time-dependent IPCW AUC.

\begin{table}[ht]
\caption{Cohort characteristics and prediction target distributions.
$\dagger$~patients with complete labels for all four classification targets.
Ext-2: data collection in progress.}
\label{tab:cohort}
\centering
\small
\begin{tabular}{lccc}
\toprule
\textbf{Characteristic}
  & \textbf{Internal} ($n=301$)
  & \textbf{Ext-1} ($n=41$)
  & \textbf{Ext-2} (TBD) \\
\midrule
Enrolment years & 2015--2019 & $>$2019 & — \\
B-ALL & 163 (54.2\%) & [X] & — \\
Patients with complete labels$^\dagger$ & 301 & 29 & — \\
\midrule
\multicolumn{4}{l}{\textit{Classification targets}} \\
Final Risk: LR & 159 (52.8\%) & [X] & — \\
Final Risk: IR/HR & 142 (47.2\%) & [X] & — \\
Relapse: Yes & 44 (14.6\%) & [X] & — \\
D19MRD $\geq 1\%$ & [X] ([X]\%) & [X] & — \\
D46MRD $\geq 0.01\%$ & [X] ([X]\%) & $\leq$2 & — \\
\midrule
\multicolumn{4}{l}{\textit{Survival endpoints}} \\
EFS events & 75 (24.9\%) & [X] & — \\
OS events & 24 (8.0\%) & [X] & — \\
\midrule
\multicolumn{4}{l}{\textit{Multimodal data}} \\
Bone marrow images & [X] (XX\%) & 41 & — \\
Morphology features & 287 (95.3\%) & 19 & — \\
\bottomrule
\end{tabular}
\end{table}

%%--------------------------------------------------------------%%
\subsection{Phase~1: LLM zero-shot performance across six targets}
\label{sec:phase1}
%%--------------------------------------------------------------%%

We evaluated four LLMs in a zero-shot setting across all twelve feature
levels and all six prediction targets (Table~\ref{tab:summary}).
Models included Gemini-3.1-Flash-Lite (Google AI API) and three
open-weight models deployed locally via vLLM: Gemma-4-31B-it,
InternVL3.5-30B-A3B, and Qwen3.6-35B-A3B.
For classification tasks (targets 1--4), each patient's feature vector
was serialised into natural language and the model output a categorical
label; we report AUROC.
For survival targets (5--6), the model output a continuous risk score
(0--1) evaluated by concordance index (C-index).

\textbf{Classification performance on the internal cohort.}
Final Risk was the easiest target for all models.
At Day~0 features ($\mathcal{F}_1$: WBC and age), Gemini and Gemma
already achieved AUROC~$\approx 0.65$--$0.70$, reflecting the known
prognostic importance of WBC count.
Performance increased incrementally through $\mathcal{F}_7$
(cytogenetics panel: Gemini~0.834, Gemma~0.891, InternVL~0.848,
Qwen~0.882), and reached its peak with MRD values included
($\mathcal{F}_8$/$\mathcal{F}_{10}$), with best-feature-level AUROCs
of 0.868 (Gemini), 0.900 (Gemma), 0.869 (InternVL), and
0.899 (Qwen).
Relapse prediction was substantially harder: best AUROCs were
0.658 (Gemini), 0.643 (InternVL), 0.586 (Gemma), and 0.580 (Qwen)
across all feature levels, with no consistent improvement as features
accumulated.
D19MRD and D46MRD showed moderate performance (AUROC~0.56--0.68),
suggesting that LLMs can partially infer early treatment response from
immunophenotypic and cytogenetic profiles but that these tasks remain
harder without direct MRD measurement.

\textbf{Survival performance on the internal cohort.}
For EFS, best-level C-indices were 0.716 (Gemma,~$\mathcal{F}_8$),
0.705 (Gemini,~$\mathcal{F}_9$), 0.694 (Qwen,~$\mathcal{F}_9$), and
0.682 (InternVL,~$\mathcal{F}_9$).
The MRD-only feature set $\mathcal{F}_9$ was consistently near-optimal
for EFS, confirming that residual disease burden encodes strong
prognostic information accessible to LLMs without labelled training
data.
For OS—with only 24 events—C-indices ranged more widely:
InternVL reached 0.772 ($\mathcal{F}_{11}$), Gemini 0.736
($\mathcal{F}_{11}$), Gemma 0.718 ($\mathcal{F}_8$), and Qwen 0.718
($\mathcal{F}_8$).
Log-rank tests were significant ($p < 0.05$) at $\mathcal{F}_8$ and
$\mathcal{F}_9$ for EFS for all four LLMs.

\textbf{External Site~1 validation.}
LLM zero-shot performance generalised well to External Site~1 for
Final Risk (AUROC~0.810--0.925) and Relapse (AUROC~0.630--0.717),
with Gemini achieving the highest Final Risk AUROC (0.925,
$\mathcal{F}_8$) and InternVL achieving the highest OS C-index
(0.859, $\mathcal{F}_8$).
D46MRD AUROCs were unstable due to very few positive cases at the
external site ($n_\text{pos} \leq 2$, results marked $\dagger$).
Survival prediction transferred robustly: EFS C-indices of
0.645--0.715 and OS C-indices of 0.806--0.859 represent generalisation
without any site-specific retraining.

%%--------------------------------------------------------------%%
\subsection{Bone marrow morphology pipeline}
\label{sec:morph_results}
%%--------------------------------------------------------------%%

Bone marrow smear images were processed through a four-stage automated
pipeline: SAM-based cell segmentation, morphological feature extraction
($\sim$37 features per cell), BiomedCLIP zero-shot cell-type clustering
($k=7$), and structured text report generation.
Comparing $\mathcal{F}_{4}$ (raw images) with $\mathcal{F}_{4.1}$
(blast-only morphology text report), the text-report approach yielded
higher Relapse AUROC for all four models (+0.017~Gemini, +0.035~Gemma,
+0.022~InternVL, +0.036~Qwen).
Final Risk AUROC remained comparable or slightly improved, and the
all-cell-type report ($\mathcal{F}_{4.2}$) improved Final Risk for
Gemma (0.854 vs 0.840 for raw images).
These results demonstrate that domain-specific cell-level
characterisation provides predictive signal complementary to global
image appearance, and that converting images to structured text
enables LLMs to extract clinically relevant morphological information.

%%--------------------------------------------------------------%%
\subsection{Phase~2: Temporal validation and classification+survival
results}
\label{sec:phase2}
%%--------------------------------------------------------------%%

Phase~2 evaluates DT and RF models on three temporal splits
($\mathcal{T}_1=2016$: $n_\text{train}=99$, $n_\text{test}=201$;
$\mathcal{T}_2=2017$: $n_\text{train}=186$, $n_\text{test}=114$;
$\mathcal{T}_3=2018$: $n_\text{train}=283$, $n_\text{test}=17$)
and a random survival forest (RSF) on the same splits for the EFS
and OS targets (Table~\ref{tab:summary}).

\textbf{Classification on the internal cohort (Split~2).}
Final Risk DT AUROC rose from 0.810 at $\mathcal{F}_1$ to
0.912 at $\mathcal{F}_7$, with the largest single gain at $\mathcal{F}_6$
(+0.127, introduction of Initial Risk and Day~5 blast percentage).
Including MRD at $\mathcal{F}_8$ further raised DT AUROC to
\textbf{0.949} (RF:~0.998), the highest classification result across
all conditions.
D19MRD achieved DT AUROC~0.743 and RF AUROC~0.807 at $\mathcal{F}_{4.1}$
(morphology-augmented); D46MRD reached DT~0.671 and RF~0.636 at
$\mathcal{F}_5$.
Relapse prediction remained difficult for Phase~2 as for Phase~1
(best DT AUROC~0.529, RF~0.630), mirroring the broader literature
on relapse biology~\cite{[CITE:Martinez2025]}.

\textbf{Survival on the internal cohort (Split~2).}
RSF at $\mathcal{F}_8$ achieved C-index~\textbf{0.746} and
time-AUC at 24~months~0.775 for EFS (log-rank $p=0.010$),
outperforming the best LLM zero-shot EFS C-index (0.716) obtained
without any training data.
For OS, RSF C-index was 0.629 ($p=0.489$; not significant), reflecting
the low event rate (8.0\%) and limited statistical power.

\textbf{External Site~1 (OOD) evaluation.}
Using all 301 internal patients as training and 41 external patients
as the test set, RF replicated its high Final Risk performance
(AUROC~0.991 at $\mathcal{F}_8$) and DT achieved 0.968, confirming
that models trained on the internal temporal splits generalise to
patients from an independent institution enrolled after the training
period.
D19MRD RF AUROC was 0.635 (external) vs 0.807 (internal Split~2),
suggesting modest feature-distribution shift.
D46MRD RF AUROC was 1.000 on External Site~1; however, this result
is unstable due to very few positive cases ($n_\text{pos}=1$, marked
$\dagger$ in Table~\ref{tab:summary}).
Relapse RF AUROC was 0.528, consistent with the difficulty observed
internally.
For survival, RSF EFS C-index was 0.608 (not significant on Ext-1),
while OS C-index reached 0.762 with log-rank $p=0.031$
($\mathcal{F}_{12.1}$; morphology-augmented).

%%--------------------------------------------------------------%%
\subsection{Cosine RAG belief revision}
\label{sec:rag}
%%--------------------------------------------------------------%%

For each test patient $i$, the $K=5$ most similar training cases were
retrieved by cosine similarity in the standardised feature space and
used to compute a similarity-weighted RAG evidence probability
$p_\text{rag}$ (Eq.~\ref{eq:prag}).
The revised probability is
$p_\text{rev} = \alpha\,p_0 + (1-\alpha)\,p_\text{rag}$, where
$\alpha = 1-(1-\alpha_\text{base})\sqrt{\bar{s}}$ and
$\alpha_\text{base}=0.6$.

\textbf{Internal cohort.}
Table~\ref{tab:belief_revision} reports AUROC before and after belief
revision.
Belief revision improved performance on Final Risk (Split~2:
$+0.002$), D19MRD (Split~2: $+0.012$; Split~3: $+0.024$), and most
substantially on D46MRD (Split~2: $+0.046$; Split~3: $0.655 \to 0.881$,
$\Delta=+0.226$).
The large D46MRD gain on Split~3 reflects high retrieval quality: with
283 training cases, the five nearest neighbours share strong
phenotypic similarity (mean cosine $\bar{s} \approx 0.91$) and
concordant D46MRD outcomes.
Belief revision degraded Relapse performance (Split~2:
$0.630 \to 0.588$, $\Delta=-0.042$): because relapse prediction is
intrinsically noisy, retrieved neighbours have mixed outcomes, and
RAG evidence acts as a noise source rather than a corrective signal.

\textbf{External Site~1.}
On external patients, belief revision improved Final Risk
($0.991 \to 1.000$, $+0.009$).
D19MRD AUROCs degraded ($-0.064$) because the external set has very
few D19MRD-positive patients ($n_\text{pos}=4$), making retrieved
neighbour outcomes uninformative.
D46MRD belief revision could not be computed ($n_\text{pos}=1$).
These results confirm the internal boundary condition: belief revision
benefits tasks where the feature space clusters meaningfully by label,
but is counter-productive when class representation is too sparse to
yield informative neighbours.

\textbf{Survival belief revision.}
For EFS and OS, belief revision was applied to normalised RSF risk
scores using event indicators as binary outcomes.
Effects were near-neutral for EFS (mean $\Delta$C-index $= -0.003$
across splits) and slightly negative for OS
(mean $\Delta$C-index $= -0.007$), where the low event rate
(8\%) renders most retrieved neighbours censored and the RAG evidence
is uninformative.

\begin{table}[ht]
\caption{Cosine RAG belief revision results.
$p_\text{RF}$: original RF AUROC;
$p_\text{rev}$: AUROC after belief revision ($K=5$,
$\alpha_\text{base}=0.6$).
Best experiment per task shown.
$\dagger$~insufficient positives; belief revision not computed.}
\label{tab:belief_revision}
\centering
\small
\begin{tabular}{llccc|ccc}
\toprule
& & \multicolumn{3}{c|}{\textbf{Internal}}
  & \multicolumn{3}{c}{\textbf{External Site~1}} \\
\textbf{Task} & \textbf{Split}
  & RF AUC & Rev AUC & $\Delta$
  & RF AUC & Rev AUC & $\Delta$ \\
\midrule
Final Risk & Split~2 & 0.998 & 0.999 & $+0.002$ & — & — & — \\
Final Risk & OOD     & — & — & — & 0.991 & 1.000 & $+0.009$ \\
D19MRD     & Split~2 & 0.767 & 0.779 & $+0.012$ & — & — & — \\
D19MRD     & Split~3 & 0.833 & 0.857 & $+0.024$ & — & — & — \\
D19MRD     & OOD     & — & — & — & 0.335 & 0.270 & $-0.064$ \\
D46MRD     & Split~2 & 0.533 & 0.579 & $+0.046$ & — & — & — \\
D46MRD     & Split~3 & 0.655 & \textbf{0.881} & $\mathbf{+0.226}$ & — & — & — \\
D46MRD     & OOD     & — & — & — & \multicolumn{3}{c}{$\dagger$ ($n_\text{pos}=1$)} \\
Relapse    & Split~2 & 0.630 & 0.588 & $-0.042$ & — & — & — \\
Relapse    & OOD     & — & — & — & 0.267 & 0.283 & $+0.016$ \\
\bottomrule
\end{tabular}
\end{table}

%%--------------------------------------------------------------%%
\subsection{LLM clinical narratives and qualitative evaluation}
\label{sec:narratives}
%%--------------------------------------------------------------%%

After belief revision, an LLM prompt was constructed for each test
patient containing: (i) the patient's serialised feature vector;
(ii) the RF-predicted label, revised probability $p_\text{rev}$, and
the human-readable DT decision path from root to leaf;
(iii) the six patient-salient features ranked by
$\text{importance}_j \times |z_{ij}|$;
and (iv) the structured profiles and outcomes of the top-5 retrieved
historical cases.
Gemini-3.1-Flash-Lite produced 3--5-sentence narratives for each
prediction; all narratives from the $\mathcal{F}_8$/Split~3 experiment
($n=17$ test patients) are available for physician review.

Two representative cases from the $\mathcal{F}_8$/Split~3 experiment
illustrate the range of narrative quality (Fig.~\ref{fig:rag}).
In Case~1 (correctly classified IR, persistent MRD), the top-5
retrieved neighbours all share elevated MRD and cytogenetic
abnormalities, and the narrative coherently attributes the high-risk
classification to MRD non-clearance corroborated by historical
precedent.
In Case~2 (borderline case), the narrative highlights conflicting
signals between MRD clearance and cytogenetic risk, suggesting
next-generation sequencing as a way to resolve ambiguity.

Quantitative physician evaluation of narrative quality is planned
and will follow the protocol in Table~\ref{tab:narrative_eval}.
Evaluation will cover: factual consistency with the patient's features,
relevance of historical case comparisons, clinical plausibility of the
interpretation, and overall helpfulness for clinical decision-making.

\begin{table}[ht]
\caption{Planned physician evaluation of LLM clinical narratives.
Each narrative will be assessed on four dimensions by two independent
clinicians using a 5-point Likert scale.
Results to be filled following structured review.}
\label{tab:narrative_eval}
\centering
\small
\begin{tabular}{lccccc}
\toprule
\textbf{Case ID}
  & \textbf{Factual accuracy}
  & \textbf{Retrieval relevance}
  & \textbf{Clinical plausibility}
  & \textbf{Decision helpfulness}
  & \textbf{Overall} \\
  & (1--5) & (1--5) & (1--5) & (1--5) & (1--5) \\
\midrule
Case 001 & — & — & — & — & — \\
Case 002 & — & — & — & — & — \\
Case 003 & — & — & — & — & — \\
$\vdots$ & $\vdots$ & $\vdots$ & $\vdots$ & $\vdots$ & $\vdots$ \\
Case 017 & — & — & — & — & — \\
\midrule
Mean (SD) & — & — & — & — & — \\
Inter-rater $\kappa$ & — & — & — & — & — \\
\bottomrule
\end{tabular}
\end{table}

%%==============================================================%%
\section{Discussion}\label{sec:discussion}
%%==============================================================%%

This study presents a temporally-structured multimodal framework
evaluated jointly across six ALL prediction targets—four classification
tasks and two survival endpoints—and validated at an independent
external centre.
Several findings warrant discussion.

\textbf{MRD as the dominant signal across all targets.}
The largest single performance gain in Phase~2—DT AUROC for Final Risk
from 0.912 ($\mathcal{F}_7$, cytogenetics) to 0.949 ($\mathcal{F}_8$,
plus D19/D46 MRD)—quantitatively confirms the central role of MRD
in clinical risk stratification~\cite{[CITE:Conter2010]}.
The same pattern holds for LLM EFS C-index, which peaked at feature
levels including MRD (Gemma: 0.716 at $\mathcal{F}_8$), and for
Phase~2 RSF EFS (C-index~0.746, $p=0.010$).
These convergent results from four LLMs, decision trees, and a survival
forest provide multi-method evidence that residual disease burden at
Day~46 is the strongest prognostic signal accessible before relapse.

\textbf{LLM zero-shot as a competitive baseline.}
Phase~1 Final Risk AUROCs of 0.868--0.900 without labelled training
data are competitive with Phase~2 DT at $\mathcal{F}_7$ (0.912) and
only modestly below the Phase~2 DT peak (0.949 at $\mathcal{F}_8$).
This indicates that clinical knowledge encoded during LLM pretraining
includes a meaningful understanding of paediatric ALL risk factors, and
that LLMs may serve as strong priors when historical case databases are
small or absent—for example, at the time a new treatment protocol is
introduced.
On external Site~1, LLM Final Risk AUROCs of 0.810--0.925 surpassed
the internal zero-shot values for some models, likely due to differences
in class distribution (larger IR/HR proportion in the external cohort)
rather than genuine improvement in discriminative ability.

\textbf{External validation: generalisation and limitations.}
Phase~2 RF Final Risk AUROC of 0.991 on External Site~1 (versus
0.998 internally) demonstrates strong cross-site generalisation.
However, D19MRD AUROC dropped from 0.807 (internal) to 0.635
(external), indicating feature-distribution shift that partly limits
generalisation for this target.
The very low positive rate for D46MRD at the external site ($n=1$)
precluded belief revision and produced unstable AUC estimates,
highlighting the need for larger external cohorts with balanced
outcome distributions.
External Site~2 data, once collected, will provide a more comprehensive
test of cross-site robustness.

\textbf{Integrating six targets in a unified framework.}
Treating EFS and OS as targets~5--6 of the same feature hierarchy—rather
than studying them in a separate survival analysis—reveals a coherent
picture: the same feature level ($\mathcal{F}_8$) that maximises
Final Risk classification also maximises EFS survival discrimination
(RSF C-index~0.746), while Relapse and OS remain challenging regardless
of feature set.
This joint view avoids the selective reporting that can arise when
survival and classification tasks are presented in separate papers or
using different patient subsets.

\textbf{Boundary conditions for RAG belief revision.}
Belief revision delivers quantitative gains specifically when the feature
space clusters meaningfully by label.
For D46MRD at Split~3 (mean similarity $\bar{s}\approx 0.91$), the
$\Delta+0.226$ AUROC gain is substantial and clinically relevant.
For Relapse ($\Delta-0.042$ internally, $+0.016$ externally with very
low base AUROC), and for OS (mean $\Delta$C-index $=-0.007$), retrieved
neighbours are too heterogeneous to provide reliable evidence.
A practical deployment criterion would be to apply belief revision
only when inter-neighbour outcome agreement exceeds a calibrated
threshold—for example, Fleiss $\kappa > 0.4$ over the top-$K$
retrieved labels.

\textbf{Relapse remains an open problem.}
Relapse prediction was intractable for all Phase~1 and Phase~2 models
(best DT AUROC~0.529, LLM best~0.658), consistent with the broader
literature~\cite{[CITE:Martinez2025],[CITE:Good2018]}.
Our feature set does not include genomic or epigenomic data (e.g.,
copy-number variation, DNA methylation), which have shown modest
incremental utility~\cite{[CITE:Mosquera2024]}.
Future multi-omics integration within the temporal feature hierarchy
may improve relapse prediction.

\textbf{Limitations.}
Both cohorts are single-institution retrospective studies; multi-centre
prospective validation is required before clinical deployment.
The internal cohort ($n=301$) limits statistical power for OS
(24 events) and Split~3 ($n_\text{test}=17$).
Phase~1 LLMs were evaluated zero-shot; fine-tuning on labelled clinical
data may yield further improvements.
Physician evaluation of LLM narratives (Table~\ref{tab:narrative_eval})
is ongoing and will determine whether the narratives are clinically
actionable beyond proof-of-concept case studies.

%%==============================================================%%
\section{Methods}\label{sec:methods}
%%==============================================================%%

\subsection{Datasets and ethics}\label{sec:dataset}

The internal training cohort comprised $n=301$ patients diagnosed with
ALL at [HOSPITAL NAME] between 2015 and 2019.
All patients were treated with risk-stratified chemotherapy protocols
consistent with contemporary Chinese paediatric oncology guidelines.
Ethics approval was obtained from [ETHICS NUMBER]; all data were
de-identified prior to analysis.
Each patient record includes structured clinical data (demographics,
haematological indices, immunophenotype, 51-marker flow cytometry panel,
cytogenetics, MRD values), bone marrow aspirate microscopy images
(Giemsa-stained, $\sim$20 fields per patient), and diagnostic
pathology reports.

External Site~1 comprised 41 patients from Lianyungang First People's
Hospital diagnosed after 2019; ethics approval: [ETHICS NUMBER Ext-1].
External Site~2 data collection is in progress ([ETHICS NUMBER Ext-2]).
No external data were used in any model training.

\subsection{Temporally-structured feature hierarchy}\label{sec:features}

Feature sets form a strictly nested sequence
$\mathcal{F}_1 \subseteq \cdots \subseteq \mathcal{F}_{12}$,
where each level contains exactly the features available at the
corresponding clinical timepoint.
Boundaries: $\mathcal{F}_1$~(Day~0): WBC count and age;
$\mathcal{F}_2$~(Day~1): adds blast percentage;
$\mathcal{F}_3$~(Day~1): adds immunophenotype and 51 flow cytometry markers;
$\mathcal{F}_6$~(Day~5): adds CNS status, Day~5 blast, and Initial Risk;
$\mathcal{F}_7$~(Day~7--10): adds Fusion gene, Karyotype, FISH;
$\mathcal{F}_8$~(Day~46): adds D19MRD and D46MRD.
Image-based variants ($\mathcal{F}_4$, $\mathcal{F}_{4.1}$, $\mathcal{F}_{10}$,
$\mathcal{F}_{12.1}$) incorporate bone marrow morphology.
All six prediction targets are evaluated on the same feature hierarchy.

For Phase~2, categorical features were one-hot encoded and numerical
features z-score normalised using training-split statistics only:
\begin{equation}
z_{ij}^{(\text{test})} =
\frac{x_{ij}^{(\text{test})} - \mu_j^{(\text{train})}}
     {\sigma_j^{(\text{train})}},
\label{eq:zscore}
\end{equation}
with missing values imputed by training-split medians.

\subsection{Bone marrow morphology pipeline}\label{sec:morph}

\textbf{Stage~1: Cell segmentation.}
SAM~3.1~\cite{[CITE:SAM]} was applied with the text prompt
\textit{``complete dark-stained nucleated cell''} after colour-based
pre-filtering (HSV hue 120--160°) to restrict sampling to
haematoxylin-stained regions.

\textbf{Stage~2: Morphological feature extraction.}
For each retained cell, $\sim$37 features were computed: geometric
descriptors (area, perimeter, circularity, eccentricity, aspect ratio),
textural descriptors (GLCM contrast and correlation), nucleus-to-cytoplasm
(N/C) ratio, nucleoli count, and vacuole percentage.

\textbf{Stage~3: Cell-type clustering.}
BiomedCLIP~\cite{[CITE:BiomedCLIP]} was applied in a zero-shot
clustering mode with $k=7$ prototype classes (blast, granulocyte,
megakaryocyte, monocyte, plasma cell, red cell, other), with Hungarian
algorithm assignment.

\textbf{Stage~4: Report generation.}
Per-patient statistics were aggregated and serialised into a structured
natural-language report supplied to LLM prompts.

\subsection{Phase~1: LLM zero-shot inference}\label{sec:llm_inference}

For classification tasks, each patient's features were serialised into a
key--value list and presented with a task-specific system prompt
requesting a bracketed label (e.g., \texttt{[[IR]]}/\texttt{[[LR]]}).
For survival risk scoring, the same serialised features were presented
with a prompt requesting a continuous event probability (0.0--1.0).
All experiments used temperature~$=0$ (greedy decoding).
Per-patient outputs were cached; the cohort was scored for every
combination of feature level ($\times 12$), task ($\times 6$), and
model ($\times 4$).

\subsection{Phase~2: Decision tree and random forest}\label{sec:dt}

A decision tree (DT; \texttt{max\_depth}=4, \texttt{class\_weight}=%
`balanced', \texttt{criterion}=`gini', \texttt{random\_state}=42)
and a random forest (RF; $n_\text{estimators}=100$,
\texttt{class\_weight}=`balanced') were trained for each feature level
and temporal split.
Three temporal splits defined by enrolment year threshold $\mathcal{T}_s$
(Eq.~\ref{eq:splits}) ensure strict temporal non-leakage.
\begin{equation}
\mathcal{D}_s^{\text{train}} = \{i : \tau_i \leq \mathcal{T}_s\},
\quad
\mathcal{D}_s^{\text{test}} = \{i : \tau_i > \mathcal{T}_s\}.
\label{eq:splits}
\end{equation}
Primary metric: AUROC; secondary: F1-macro and balanced accuracy.

\subsection{Phase~2: Random survival forest}\label{sec:rsf}

A random survival forest (RSF;~\cite{[CITE:Ishwaran2008]};
$n_\text{estimators}=100$, \texttt{random\_state}=42) was trained on
the same temporal splits as the classification experiments.
EFS was defined as time from diagnosis to first of: relapse, induction
failure, or death; OS was time to death.
Censoring was at last follow-up.
Evaluation used Harrell's C-index (Eq.~\ref{eq:cindex}), IPCW
time-dependent AUC (Eq.~\ref{eq:itauc}) at $t^*\in\{12,24,36\}$
months, and the log-rank test on Kaplan-Meier curves dichotomised at
the median risk score.

\subsection{External OOD evaluation}\label{sec:ood}

For external validation, all $n=301$ internal patients served as the
training set and external patients served as the test set; no
site-specific fine-tuning was performed.
DT and RF were trained on the full internal set using the same
preprocessing pipeline (Eq.~\ref{eq:zscore}).
RSF was trained identically.
For Phase~1 LLM evaluation, no training was involved; the same
zero-shot prompts were applied to external patient feature vectors.

\subsection{RAG belief revision and LLM narration}\label{sec:rag_methods}

For test patient $i$, the $K=5$ most similar training cases were
retrieved by cosine similarity in the standardised feature space:
\begin{equation}
\mathcal{N}(i) = \underset{j \in \mathcal{D}^{\text{train}}}
  {\arg\text{top-}K}\
  \frac{\mathbf{z}_i \cdot \mathbf{z}_j}
       {\|\mathbf{z}_i\|_2 \|\mathbf{z}_j\|_2}, \quad K = 5.
\label{eq:rag}
\end{equation}
The RAG evidence probability is
\begin{equation}
p_{\text{rag}} =
  \frac{\sum_{j \in \mathcal{N}(i)} s_j\, y_j}
       {\sum_{j \in \mathcal{N}(i)} s_j},
\label{eq:prag}
\end{equation}
where $s_j$ is the cosine similarity to neighbour $j$ and
$y_j \in \{0,1\}$ is the binary task outcome (or event indicator
for survival).
The revised probability is
\begin{equation}
p_{\text{rev}} = \alpha\, p_0 + (1-\alpha)\, p_{\text{rag}},
\qquad
\alpha = 1 - (1-\alpha_{\text{base}})\sqrt{\bar{s}},
\label{eq:prev}
\end{equation}
where $\alpha_\text{base}=0.6$ and
$\bar{s} = \frac{1}{K}\sum_{j \in \mathcal{N}(i)} s_j$.
For survival, $p_0$ is the min-max normalised RSF risk score.

After belief revision, an LLM prompt containing the patient's features,
DT decision path, top-6 salient features (ranked by
$\text{importance}_j \times |z_{ij}|$), and top-5 retrieved case
profiles was submitted to Gemini-3.1-Flash-Lite to generate a 3--5
sentence clinical narrative.

\subsection{Survival metrics}\label{sec:surv_methods}

Harrell's concordance index:
\begin{equation}
\hat{C} =
\frac{
  \sum_{i \neq j}
  \mathbf{1}[T_i < T_j]\,
  \mathbf{1}[\delta_i = 1]\,
  \mathbf{1}[\hat{r}_i > \hat{r}_j]
}{
  \sum_{i \neq j}
  \mathbf{1}[T_i < T_j]\,
  \mathbf{1}[\delta_i = 1]
}.
\label{eq:cindex}
\end{equation}
IPCW time-dependent AUC:
\begin{equation}
\widehat{\text{AUC}}(t^*) =
\frac{
  \sum_{i:\,T_i \leq t^*,\,\delta_i=1}
  \sum_{j:\,T_j > t^*}
  \mathbf{1}[\hat{r}_i > \hat{r}_j]\,\hat{w}_i
}{
  \sum_{i:\,T_i \leq t^*,\,\delta_i=1}
  \sum_{j:\,T_j > t^*}
  \hat{w}_i
},
\label{eq:itauc}
\end{equation}
where $\hat{w}_i = 1/\hat{G}(T_i)^2$ and $\hat{G}(\cdot)$ is the
Kaplan-Meier censoring estimator~\cite{[CITE:Uno2011]}.

\subsection{Statistical analysis}

AUROC was computed with \texttt{scikit-learn}~\cite{[CITE:sklearn]};
C-index and IPCW AUC with
\texttt{scikit-survival}~\cite{[CITE:sksurv]}.
Bootstrap 95\% confidence intervals (1000 resamples) are reported for
Split~3 ($n_\text{test}=17$).
No multiple-testing correction was applied given the exploratory nature
of this study; all reported $p$-values are two-sided.

%%==============================================================%%
\section*{Data availability}
%%==============================================================%%

The clinical dataset contains patient health information and cannot be
made publicly available due to privacy regulations.
De-identified data may be available through a data access agreement
with [HOSPITAL NAME]; requests to the corresponding author.

%%==============================================================%%
\section*{Code availability}
%%==============================================================%%

Analysis code, LLM prompts, and the bone marrow morphology pipeline
will be made available at
\url{https://github.com/YanbeiJiang/ALL_prediction_NCS}
upon acceptance.

%%==============================================================%%
\section*{Acknowledgements}
%%==============================================================%%

[To be completed — funding sources, clinical collaborators]

%%==============================================================%%
\section*{Author contributions}
%%==============================================================%%

[To be completed]

%%==============================================================%%
\section*{Competing interests}
%%==============================================================%%

The authors declare no competing interests.

%%==============================================================%%
%% Main tables
%%==============================================================%%

\begin{table}[ht]
\caption{Feature set design aligned with the clinical data-acquisition
timeline. All six prediction targets are evaluated at every feature
level; $\dagger$~at $\mathcal{F}_8$/$\mathcal{F}_9$, D19MRD and
D46MRD are used as \emph{input features}, so only Final Risk and
Relapse are predicted. EFS/OS are evaluated at all levels using
the same feature vector as classification tasks.}
\label{tab:expdesign}
\centering
\small
\begin{tabular}{lllc}
\toprule
\textbf{Level} & \textbf{Timepoint} & \textbf{Cumulative features added} & \textbf{Class. targets} \\
\midrule
$\mathcal{F}_1$     & Day 0     & WBC count, Age & 4 \\
$\mathcal{F}_2$     & Day 1     & + Blast\% & 4 \\
$\mathcal{F}_3$     & Day 1     & + Immunophenotype + 51 flow cytometry markers & 4 \\
$\mathcal{F}_{4}$   & Day 1     & Same as $\mathcal{F}_3$ + bone marrow images & 4 \\
$\mathcal{F}_{4.1}$ & Day 1     & Same as $\mathcal{F}_3$ + blast morphology report & 4 \\
$\mathcal{F}_5$     & Day 1     & $\mathcal{F}_3$ + non-risk clinical factors & 4 \\
$\mathcal{F}_6$     & Day 5     & $\mathcal{F}_3$ + CNS status + Day~5 blast\% + Initial Risk & 4 \\
$\mathcal{F}_7$     & Day 7--10 & $\mathcal{F}_6$ + Fusion gene + Karyotype + FISH & 4 \\
$\mathcal{F}_8$     & Day 46    & $\mathcal{F}_7$ + D19MRD + D46MRD (as inputs) & 2$^\dagger$ \\
$\mathcal{F}_9$     & Day 46    & D19MRD + D46MRD only & 2$^\dagger$ \\
$\mathcal{F}_{10}$  & Day 7--10 & $\mathcal{F}_7$ + diagnostic pathology report & 4 \\
$\mathcal{F}_{12.1}$& Day 7--10 & $\mathcal{F}_7$ + blast morphology report & 4 \\
\bottomrule
\multicolumn{4}{l}{Survival targets (EFS, OS) evaluated at all levels using the same feature vector.}
\end{tabular}
\end{table}

\begin{table}[ht]
\caption{Summary of best results across six prediction targets, four
LLMs (Phase~1), and Phase~2 models (DT/RF for classification;
RSF for EFS/OS).
Internal: best-feature-level result; all four models shown.
External~1: OOD results using all 301 internal patients as training.
Classification metric: AUROC. Survival metric: C-index.
$\dagger$~fewer than 3 positives in external cohort; AUC/C-index
unstable.
Phase~2 internal: Split~2 ($n_\text{train}=186$, $n_\text{test}=114$).
Ext-2: pending.}
\label{tab:summary}
\centering
\footnotesize
\setlength{\tabcolsep}{4pt}
\begin{tabular}{llccccc|cc|cc}
\toprule
& & \multicolumn{4}{c}{\textbf{Phase~1 LLM (best feature level)}}
  & & \multicolumn{2}{c|}{\textbf{Phase~2 (Int.)}}
  & \multicolumn{2}{c}{\textbf{Phase~2 (Ext-1)}} \\
\cmidrule(lr){3-6} \cmidrule(lr){8-9} \cmidrule(lr){10-11}
\textbf{Task} & \textbf{Metric}
  & Gemini & Gemma & InternVL & Qwen
  & \textit{(Ext-1)}
  & DT & RF/RSF & DT & RF/RSF \\
\midrule
\multicolumn{11}{l}{\textit{Classification tasks (AUROC)}} \\
Final Risk & Int. & 0.868 & 0.900 & 0.869 & 0.899 & & 0.949 & 0.998 & — & — \\
           & Ext-1 & 0.925 & 0.894 & 0.810 & 0.850 & & — & — & 0.968 & 0.991 \\
D19MRD     & Int. & 0.679 & 0.683 & 0.616 & 0.635 & & 0.743 & 0.807 & — & — \\
           & Ext-1 & 0.550 & 0.465 & 0.521 & 0.604 & & — & — & 0.628 & 0.635 \\
D46MRD     & Int. & 0.661 & 0.589 & 0.575 & 0.564 & & 0.671 & 0.636 & — & — \\
           & Ext-1 & \multicolumn{4}{c}{$\leq$0.554\,$^\dagger$}  & & — & — & 0.975$^\dagger$ & 1.000$^\dagger$ \\
Relapse    & Int. & 0.658 & 0.586 & 0.643 & 0.580 & & 0.529 & 0.630 & — & — \\
           & Ext-1 & 0.717 & 0.656 & 0.699 & 0.630 & & — & — & 0.583 & 0.528 \\
\midrule
\multicolumn{11}{l}{\textit{Survival tasks (C-index)}} \\
EFS        & Int. & 0.705 & 0.716 & 0.682 & 0.694 & & — & 0.746 & — & — \\
           & Ext-1 & 0.715 & 0.669 & 0.709 & 0.645 & & — & — & — & 0.608 \\
OS         & Int. & 0.736 & 0.718 & 0.772 & 0.718 & & — & 0.629 & — & — \\
           & Ext-1 & 0.810 & 0.821 & 0.859 & 0.806 & & — & — & — & 0.762 \\
\bottomrule
\end{tabular}
\end{table}

%%==============================================================%%
%% Figure captions
%%==============================================================%%

\begin{figure}[ht]
\caption{\textbf{Overview of the temporally-structured multimodal
prediction framework.}
\textbf{a,}~The three parallel timelines: feature acquisition
(Day~0--46+, $\mathcal{F}_1$–$\mathcal{F}_{12}$), prediction targets
(six tasks: Final Risk, D19MRD, D46MRD, Relapse, EFS, OS), and the
dynamic case database (patient enrolment year 2015--2019, temporal
splits).
\textbf{b,}~Phase~1 architecture: four LLMs receive serialised patient
features and output class labels or risk scores without task-specific
training.
\textbf{c,}~Phase~2 architecture: a decision tree trained on historical
cases generates predictions and decision paths; RAG retrieves the
$K=5$ most similar cases; LLM synthesises a clinical narrative.
\textbf{d,}~Evaluation scope: internal temporal splits (three splits) and
external OOD validation (two independent sites; Site~2 pending).}
\label{fig:overview}
\end{figure}

\begin{figure}[ht]
\caption{\textbf{Phase~1 feature-level ablation across six prediction
targets.}
\textbf{a,}~Final Risk AUROC vs feature level for four LLMs.
\textbf{b,}~Relapse AUROC. \textbf{c,}~D19MRD AUROC.
\textbf{d,}~D46MRD AUROC. \textbf{e,}~EFS C-index.
\textbf{f,}~OS C-index.
Vertical dashed lines indicate clinical timepoints: Day~1 (flow
cytometry), Day~5 (early response), Day~7--10 (cytogenetics), Day~46
(MRD, $\mathcal{F}_8$/$\mathcal{F}_9$).
Open symbols: internal cohort. Filled symbols: External Site~1.}
\label{fig:ablation}
\end{figure}

\begin{figure}[ht]
\caption{\textbf{Bone marrow morphology pipeline and its predictive
contribution.}
\textbf{a,}~Pipeline schematic: SAM cell segmentation
$\rightarrow$ morphological feature extraction $\rightarrow$
BiomedCLIP cell-type clustering $\rightarrow$ structured text report.
\textbf{b,}~Comparison of Relapse AUROC across three image-handling
strategies: raw images ($\mathcal{F}_4$), blast-only text report
($\mathcal{F}_{4.1}$), and all-cell-type text report
($\mathcal{F}_{4.2}$) for four LLMs.
\textbf{c,}~Same comparison for Final Risk AUROC.}
\label{fig:morph}
\end{figure}

\begin{figure}[ht]
\caption{\textbf{Phase~2 temporal learning curves and classification
results.}
\textbf{a,}~Final Risk DT AUROC across feature levels for three
temporal splits; shaded band: 95\% bootstrap CI for Split~3
($n_\text{test}=17$).
\textbf{b,}~Training-set size vs AUROC for Final Risk and Relapse
(DT and RF) at $\mathcal{F}_8$.
\textbf{c,}~External Site~1 RF AUROC comparison with internal
Split~2 RF AUROC across all four classification tasks.}
\label{fig:p2ablation}
\end{figure}

\begin{figure}[ht]
\caption{\textbf{Survival prediction: Kaplan-Meier curves and
C-index summary.}
\textbf{a,}~EFS Kaplan-Meier curves (high vs low risk, median split)
for the best Phase~1 model (Gemma, $\mathcal{F}_8$, internal cohort).
\textbf{b,}~EFS KM curves for Phase~2 RSF (Split~2, $\mathcal{F}_8$,
log-rank $p=0.010$).
\textbf{c,}~OS KM curves for Phase~2 RSF on External Site~1
(log-rank $p=0.031$).
\textbf{d,}~C-index heatmap: all LLMs $\times$ all feature levels
for EFS (internal, left; Ext-1, right).
Numbers at risk and 95\% CI shown.}
\label{fig:km}
\end{figure}

\begin{figure}[ht]
\caption{\textbf{RAG belief revision and representative case narratives.}
\textbf{a,}~AUROC before vs after belief revision across four tasks
and three splits; point size proportional to training-set size.
Filled points: internal; open points: External Site~1.
\textbf{b--c,}~Two representative cases from $\mathcal{F}_8$/Split~3.
For each case: patient feature summary, DT decision path, top-3
retrieved similar historical cases with outcomes, and LLM-generated
clinical narrative.
Case~b: correctly classified IR patient with persistent MRD and
concordant historical neighbours.
Case~c: borderline case with conflicting MRD clearance vs cytogenetic
risk signals.}
\label{fig:rag}
\end{figure}

%%==============================================================%%
%% References
%%==============================================================%%
\begin{thebibliography}{99}

\bibitem{[CITE:Hunger2015]}
Hunger, S.P. \& Mullighan, C.G.
Acute lymphoblastic leukemia in children.
\textit{N. Engl. J. Med.} \textbf{373}, 1541--1552 (2015).

\bibitem{[CITE:Bhojwani2013]}
Bhojwani, D. \& Pui, C.-H.
Relapsed childhood acute lymphoblastic leukaemia.
\textit{Lancet Oncol.} \textbf{14}, e205--e217 (2013).

\bibitem{[CITE:Pui2008]}
Pui, C.-H., Robison, L.L. \& Look, A.T.
Acute lymphoblastic leukaemia.
\textit{Lancet} \textbf{371}, 1030--1043 (2008).

\bibitem{[CITE:Good2018]}
Good, Z. et al.
Single-cell developmental classification of B cell precursor acute
lymphoblastic leukemia at diagnosis reveals predictors of relapse.
\textit{Nat. Med.} \textbf{24}, 474--483 (2018).

\bibitem{[CITE:Fitter2021]}
Fitter, S. et al.
CKLF and IL1B transcript levels at diagnosis are predictive of
relapse in children with pre-B-cell acute lymphoblastic leukaemia.
\textit{Br. J. Haematol.} \textbf{193}, 107--118 (2021).

\bibitem{[CITE:Mosquera2024]}
Mosquera Orgueira, A. et al.
Refining risk prediction in pediatric acute lymphoblastic leukemia
through DNA methylation profiling.
\textit{Clin. Epigenetics} \textbf{16}, 53 (2024).

\bibitem{[CITE:Pan2017]}
Pan, L. et al.
Machine learning applications for prediction of relapse in childhood
acute lymphoblastic leukemia.
\textit{Sci. Rep.} \textbf{7}, 7402 (2017).

\bibitem{[CITE:Martinez2025]}
Mart\'inez-Rubio, \'{A}. et al.
Computational flow cytometry immunophenotyping at diagnosis is unable
to predict relapse in childhood B-cell acute lymphoblastic leukemia.
\textit{Comput. Biol. Med.} \textbf{185}, 109831 (2025).

\bibitem{[CITE:Steyerberg2019]}
Steyerberg, E.W.
\textit{Clinical Prediction Models} 2nd edn (Springer, 2019).

\bibitem{[CITE:Phaterpekar2026]}
Phaterpekar, T. et al.
Investigating fine-tuning versus zero-shot learning for general large
language models when predicting cancer survival from initial oncology
consultation documents.
\textit{ESMO Real World Data Digit. Oncol.} (2026).
doi:10.1016/j.esmorw.2026.100703.

\bibitem{[CITE:Conter2010]}
Conter, V. et al.
Molecular response to treatment redefines all prognostic factors in
children and adolescents with B-cell precursor acute lymphoblastic
leukemia: results in 3184 patients of the AIEOP-BFM ALL 2000 study.
\textit{Blood} \textbf{115}, 3206--3214 (2010).

\bibitem{[CITE:Rudin2019]}
Rudin, C.
Stop explaining black box machine learning models for high stakes
decisions and use interpretable models instead.
\textit{Nat. Mach. Intell.} \textbf{1}, 206--215 (2019).

\bibitem{[CITE:SAM]}
Ravi, N. et al.
SAM~2: Segment anything in images and videos.
\textit{arXiv} 2408.00714 (2024).

\bibitem{[CITE:BiomedCLIP]}
Zhang, S. et al.
BiomedCLIP: a multimodal biomedical foundation model pretrained from
fifteen million scientific image-text pairs.
\textit{arXiv} 2303.00915 (2023).

\bibitem{[CITE:vLLM]}
Kwon, W. et al.
Efficient memory management for large language model serving with
PagedAttention.
\textit{Proc. ACM SOSP} (2023).

\bibitem{[CITE:Uno2011]}
Uno, H. et al.
On the C-statistics for evaluating overall adequacy of risk prediction
procedures with censored survival data.
\textit{Stat. Med.} \textbf{30}, 1105--1117 (2011).

\bibitem{[CITE:Ishwaran2008]}
Ishwaran, H. et al.
Random survival forests.
\textit{Ann. Appl. Stat.} \textbf{2}, 841--860 (2008).

\bibitem{[CITE:sklearn]}
Pedregosa, F. et al.
Scikit-learn: machine learning in Python.
\textit{J. Mach. Learn. Res.} \textbf{12}, 2825--2830 (2011).

\bibitem{[CITE:sksurv]}
P\"olsterl, S.
scikit-survival: a library for time-to-event analysis built on top
of scikit-learn.
\textit{J. Mach. Learn. Res.} \textbf{21}, 1--6 (2020).

\end{thebibliography}

\end{document}
